
\subsection*{7.5. «~Conservation du nombre~».}
Supposons maintenant \(f\) propre, donc \(X_t\) et \(X_s\) propres sur \(k(t)\) resp. \(k(s)\), de sorte qu'on dispose d'homomorphismes 2.3 :
\[
  \chi_t:K_\bullet(X_t)\longrightarrow \mathbb Z,\qquad
  \chi_s:K_\bullet(X_s)\longrightarrow \mathbb Z .
\]
Ceci dit, on a la relation
\begin{equation}\tag{7.5.1}
  \chi_s(\sigma_\bullet(x))=\chi_t(x)\qquad\text{pour } x\in K_\bullet(X_t)
\end{equation}
(«~théorème de conservation du nombre~»), en d'autres termes, on a commutativité dans le diagramme
\begin{equation}\tag{7.5.2}
\begin{tikzcd}[column sep=large,row sep=large]
K_\bullet(X_t) \arrow[r,"\sigma_\bullet"] \arrow[dr,"\chi_t"'] &
K_\bullet(X_s) \arrow[d,"\chi_s"]\\
& \mathbb Z .
\end{tikzcd}
\end{equation}
Pour prouver (7.5.1), on est encore ramené au cas où \(x=\cl_\bullet(Y_t)\), \(Y_t\) sous-schéma fermé intègre de \(X_t\), et avec les notations de la démonstration de 7.4 on est ramené à prouver que
\[
  \chi(Y_s)=\chi(Y_t),
\]
ce qui est bien connu (EGA III 7.9.4, ou Exposé III 5.7).

\subsection*{7.6.}
Il est immédiat que l'on a
\begin{equation}\tag{7.6.1}
  \sigma_\bullet\bigl(\cl_\bullet(\underline{O}_{X_t})\bigr)=\cl_\bullet(\underline{O}_{X_s}),
\end{equation}
ce qui entraîne, compte tenu que \(\sigma_\bullet\) est \(K^{\bullet}(X)\)-linéaire, que le diagramme suivant est commutatif :
\begin{equation}\tag{7.6.2}
\begin{tikzcd}[column sep=large,row sep=large]
& K^{\bullet}(X) \arrow[dl,"j^*"'] \arrow[dr,"i^*"] &\\
K^{\bullet}(X_t) \arrow[d] && K^{\bullet}(X_s) \arrow[d]\\
K_\bullet(X_t) \arrow[rr,"\sigma_\bullet"] && K_\bullet(X_s),
\end{tikzcd}
\end{equation}
